% ==========================================
% ANEXOS — ARCHIVO SEMI-ESTÁTICO
% ==========================================
% Agrega tus apéndices aquí. A diferencia de los otros
% archivos en common/, este SÍ se edita por proyecto.
%
% TIP: Los anexos se numeran con letras (A, B, C...).
%      Usa \chapter{} para cada anexo y \section{} para subdivisiones.
%
% TIP: Cada anexo que agregues aparecerá automáticamente
%      en el índice general.
% ==========================================

\appendix
\renewcommand{\appendixname}{Anexo}
\renewcommand{\appendixpagename}{Anexos}
\renewcommand{\appendixtocname}{Anexos}
\addappheadtotoc
\appendixpage

\chapter{Códigos completos}
\label{chap:codigos_completos}
% Ejemplos:


% A continuación se presentan de manera compacta los archivos que conforman el programa desarrollado de cálculo concurrente con POSIX threads, incluyendo tanto la configuración de directivas en \texttt{Makefile}, como las definiciones estructurales de datos compartidos y la lógica central en lenguaje C. 

\section{Código fuente base de hilo padre e hijos (main.c)}
\lstinputlisting[language=C, style=mystyle]{codigos/main.c}

\section{Código de configuración principal (Makefile)}
\lstinputlisting[language=make, style=mystyle]{codigos/Makefile}

% \section{Definición de librerías y estructuras (matriz\_hilos.h)}
% \lstinputlisting[language=C, style=mystyle]{codigos/matriz_hilos.h}





% En este anexo se presenta el código fuente completo de la práctica.

% \section{Makefile}
% \lstinputlisting[
% 	language=make,
% 	frame=single,
% 	basicstyle=\ttfamily\scriptsize,
% 	caption={Código fuente de \texttt{Makefile}. \propia{}},
% 	label={lst:codigo_makefile}
% ]{contexto/practica3final/Makefile}

% \section{comun.h}
% \lstinputlisting[
% 	language=C,
% 	frame=single,
% 	basicstyle=\ttfamily\scriptsize,
% 	caption={Código fuente de \texttt{comun.h}. \propia{}},
% 	label={lst:codigo_comun_h}
% ]{contexto/practica3final/comun.h}

% \section{Coordinador.c}
% \lstinputlisting[
% 	language=C,
% 	frame=single,
% 	basicstyle=\ttfamily\scriptsize,
% 	caption={Código fuente de \texttt{Coordinador.c}. \propia{}},
% 	label={lst:codigo_coordinador}
% ]{contexto/practica3final/Coordinador.c}

% \section{TerminalA.c}
% \lstinputlisting[
% 	language=C,
% 	frame=single,
% 	basicstyle=\ttfamily\scriptsize,
% 	caption={Código fuente de \texttt{TerminalA.c}. \propia{}},
% 	label={lst:codigo_terminalA}
% ]{contexto/practica3final/TerminalA.c}

% \section{TerminalB.c}
% \lstinputlisting[
% 	language=C,
% 	frame=single,
% 	basicstyle=\ttfamily\scriptsize,
% 	caption={Código fuente de \texttt{TerminalB.c}. \propia{}},
% 	label={lst:codigo_terminalB}
% ]{contexto/practica3final/TerminalB.c}

% \section{TerminalC.c}
% \lstinputlisting[
% 	language=C,
% 	frame=single,
% 	basicstyle=\ttfamily\scriptsize,
% 	caption={Código fuente de \texttt{TerminalC.c}. \propia{}},
% 	label={lst:codigo_terminalC}
% ]{contexto/practica3final/TerminalC.c}